% Modèle d'une source réelle de tension (source idéale E en série avec r) et caractéristique U = E - r I (L'essentiel, chapitre 14)
\begin{tikzpicture}[x=1cm,y=1cm]
% Modèle électrique
\begin{scope}
  \draw[otMuted!80,line width=0.5pt,dash pattern=on 2pt off 1.5pt,rounded corners=1.5mm] (-0.9,-0.35) rectangle (0.75,3.35);
  \node[pcLegende,anchor=south,text=otMuted] at (-0.08,3.38) {source réelle};
  \draw[otInk!80,line width=0.7pt] (0,0) -- (0,0.62) (0,0.78) -- (0,1.55) (0,2.45) -- (0,3.0) -- (2.6,3.0) -- (2.6,1.92) (2.6,1.08) -- (2.6,0) -- (0,0);
  \draw[otInk!80,line width=0.7pt] (-0.32,0.78) -- (0.32,0.78);
  \draw[otInk!80,line width=1.6pt] (-0.16,0.62) -- (0.16,0.62);
  \draw[otNavy,line width=0.6pt,-{Stealth[length=1.5mm]}] (-0.55,0.25) -- (-0.55,1.2) node[pos=0.5,pcLegende,anchor=east,text=otNavy] {$E$};
  \filldraw[fill=white,draw=otInk!80,line width=0.7pt] (-0.14,1.55) rectangle (0.14,2.45);
  \node[pcLegende,anchor=west] at (0.18,2.0) {$r$};
  \fill[otInk] (0,3.0) circle (0.05); \node[pcLegende,anchor=south east] at (0,3.0) {P};
  \fill[otInk] (0,0) circle (0.05); \node[pcLegende,anchor=north east] at (0,0) {N};
  \filldraw[fill=white,draw=otInk!80,line width=0.7pt] (2.46,1.08) rectangle (2.74,1.92);
  \node[pcLegende,anchor=west] at (2.78,1.5) {récepteur};
  \draw[otRed!80!black,line width=0.7pt,-{Stealth[length=1.6mm]}] (1.1,3.0) -- (1.6,3.0) node[pcLegende,anchor=south,text=otRed!80!black] {$I$};
  \draw[otNavy,line width=0.6pt,-{Stealth[length=1.5mm]}] (1.2,0.3) -- (1.2,2.7) node[pos=0.5,pcLegende,anchor=west,text=otNavy] {$U=U_\mathrm{PN}$};
\end{scope}
% Caractéristique
\begin{scope}[shift={(5.6,0)}]
  \draw[otMuted,line width=0.5pt,-{Stealth[length=2mm]}] (0,0) -- (4.4,0) node[pcLegende,anchor=north west] {$I$};
  \draw[otMuted,line width=0.5pt,-{Stealth[length=2mm]}] (0,0) -- (0,3.3) node[pcLegende,anchor=south east] {$U$};
  \draw[otNavy,line width=1.2pt] (0,2.7) -- (3.8,0);
  \node[pcLegende,anchor=east] at (0,2.7) {$E$};
  \node[pcLegende,anchor=north] at (3.8,0) {$I_\mathrm{cc}=\dfrac{E}{r}$};
  \draw[otGreenDark,line width=0.5pt] (1.0,1.99) -- (2.0,1.99) -- (2.0,1.28);
  \node[pcLegende,anchor=west,text=otGreenDark] at (2.05,1.62) {pente $-r$};
  \node[pcLegende,anchor=south west,text=otNavy] at (1.6,2.2) {$U=E-r\,I$};
\end{scope}
\end{tikzpicture}
\par\smallskip\legende{La tension aux bornes d'une source réelle diminue quand elle débite davantage : une partie de la tension « se perd » dans sa résistance interne.}
